% !TeX root = ../../Book4.tex
% 纲要标题：### 5.4 代数对象、模对象与富集
% 纲要标签：b4:sec:monoidal-algebras
% ---
% 纲要细目（写作范围）：
% #### 5.4.1 乘法与单位的对象形式
% #### 5.4.2 模对象与函子复合中的单子
% #### 5.4.3 内部复合与富集范畴

\section{代数对象、模对象与富集}
\label{b4:sec:monoidal-algebras}

普通代数的乘法是双线性映射，因而可以写成 \(A\otimes A\to A\)。其结合律比较三重乘积的两个求值过程，单位则来自基环到 \(A\) 的映射。这样写以后，运算的定义不再依赖元素记号：只要一个范畴有适当的张量与单位，就能讨论其中的代数与模。

\subsection{乘法与单位的对象形式}
\label{b4:subsec:monoidal-algebra-objects}

\begin{definition}\label{b4:def:monoidal-algebra-object}
设 \(\mathcal C\) 为幺半范畴，单位为 \(I\)。给定对象 \(A\) 及态射 \(m:A\otimes A\to A\)、\(u:I\to A\)。若满足
\begin{equation}\label{b4:eq:monoidal-algebra-laws}
\begin{aligned}
m\circ(m\otimes1_A)&=m\circ(1_A\otimes m)\circ\alpha_{A,A,A},\\
m\circ(u\otimes1_A)&=\lambda_A,\qquad
m\circ(1_A\otimes u)=\rho_A,
\end{aligned}
\end{equation}
则称 \((A,m,u)\) 为 \(\mathcal C\) 中的\term[daishu-duixiang]{代数对象}[algebra object]，也称幺半群对象。对两个代数对象，若态射 \(f:A\to B\) 满足
\[
f\circ m_A=m_B\circ(f\otimes f),\qquad f\circ u_A=u_B,
\]
则称 \(f\) 为代数对象的同态。在辫状幺半范畴中，若还满足 \(m\circ c_{A,A}=m\)，则称代数对象为交换的。
\end{definition}

第一式从 \((A\otimes A)\otimes A\) 出发：左边先乘前两个因子，再与第三个相乘；右边先用 \(\alpha\) 改括号，乘后两个因子，再与第一个相乘。它要求这两种乘法过程相同。第二行则先以 \(u\) 把单位对象变成代数中的单位，再与原输入相乘，所得映射分别等于 \(\lambda_A:I\otimes A\to A\)、\(\rho_A:A\otimes I\to A\)。由于定义域分别为 \(I\otimes A\)、\(A\otimes I\)，这里不能直接写成 \(1_A\)。这一形式可对照\cite[Definition 7.8.1, p. 141]{EtingofEtAl2015TensorCategories}；其定义所在的张量范畴带有其他有限性条件，但上述乘法公理本身只需幺半结构。

\begin{proposition}\label{b4:prop:monoidal-ordinary-algebras}
设 \(R\) 为交换含幺环。在 \((R\text{-}\mathbf{Mod},\otimes_R,R)\) 中，代数对象及其同态恰好给出含幺结合 \(R\) 代数及其保幺 \(R\) 代数同态。其交换代数对象恰好是交换 \(R\) 代数。在 \((\mathbf{Set},\times,\{*\})\) 中，代数对象及同态恰好是幺半群及幺半群同态。
\end{proposition}
\begin{proof}
给 \(R\)-模代数对象，定义 \(ab=m(a\otimes b)\)、\(1_A=u(1)\)。张量积泛性质使该乘法双线性。将\eqref{b4:eq:monoidal-algebra-laws}作用于 \((a\otimes b)\otimes c\)，得到 \((ab)c=a(bc)\)；作用于 \(1\otimes a\)、\(a\otimes1\)，得到 \(1_Aa=a=a1_A\)。又因 \(u\) 线性，\(u(r)=r1_A\)，于是
\[
u(r)a=ra=au(r),\qquad u(r)u(s)=u(rs).
\]
所以 \(u:R\to A\) 是取值于中心的保幺环同态，得到 \(R\) 代数。标量作用 \(ra\) 原先已由 \(R\) 模 \(A\) 给定；\(u(r)a=ra\) 证明沿 \(u\) 诱导的标量作用与这个原有作用相同。

反过来，含幺结合 \(R\) 代数的乘法双线性且平衡，故诱导唯一的 \(m:A\otimes_RA\to A\)。结构同态 \(u\) 线性；结合、单位公理在纯张量上成立，故给出全部所需态射等式。两构造在元素与结构映射上互逆。代数对象同态的两个条件正是保持乘法与单位，且其本身是 \(R\) 线性映射，所以恰是保幺 \(R\) 代数同态。交换约束把 \(a\otimes b\) 换成 \(b\otimes a\)，故附加条件等价于 \(ab=ba\)。

对集合，\(m\) 是二元运算，\(u\) 指定元素 \(u(*)\)。同样的等式直接成为结合律与左右单位律，同态条件也直接成为幺半群同态的条件。
\end{proof}

交换性取决于所选辫状结构。例如，在分次向量空间的符号对称结构中，对齐次元素 \(a\in A_m\)、\(b\in A_n\)，条件 \(m\circ c=m\) 写成
\[
ab=(-1)^{mn}ba.
\]
它允许两个奇次数元素的乘法带负号，因而与不分次数的普通交换性不同。对奇次数元素 \(a\)，直接得到的是 \(a^2=-a^2\)，即 \(2a^2=0\)。若 \(2\) 可逆，则由此推出 \(a^2=0\)；特征 \(2\) 时符号为正，不能这样推出平方为零。

\subsection{模对象与函子复合中的单子}
\label{b4:subsec:monoidal-module-objects}

一个代数对象要作用在另一个对象上，只须把乘法的第二个因子改成该对象，并保留同样的结合与单位要求。

\begin{definition}\label{b4:def:monoidal-module-object}
设 \((A,m,u)\) 为幺半范畴 \(\mathcal C\) 中的代数对象。给对象 \(M\) 及态射 \(a:A\otimes M\to M\)。若满足
\[
a\circ(m\otimes1_M)=a\circ(1_A\otimes a)\circ\alpha_{A,A,M},
\qquad a\circ(u\otimes1_M)=\lambda_M,
\]
则称 \((M,a)\) 为左 \(A\)-\term[mo-duixiang]{模对象}[module object]。若 \(f:M\to N\) 满足 \(f\circ a_M=a_N\circ(1_A\otimes f)\)，则称其为模对象同态。右模对象以 \(M\otimes A\to M\) 的作用定义，其结合等式为
\(a\circ(a\otimes1_A)=a\circ(1_M\otimes m)\circ\alpha_{M,A,A}\)，单位等式为 \(a\circ(1_M\otimes u)=\rho_M\)。
\end{definition}

从 \((A\otimes A)\otimes M\) 出发，作用结合式的左边先乘两个 \(A\) 因子，再作用于 \(M\)；右边改括号后，先让第二个 \(A\) 因子作用于 \(M\)，再让第一个作用。单位式要求插入代数单位后的作用恢复原输入。在交换环 \(R\) 的模范畴中，左模对象恢复通常的左 \(A\) 模。确切地，作用态射对应 \(R\) 双线性作用 \((x,v)\mapsto xv\)，两条等式成为 \((xy)v=x(yv)\)、\(1_Av=v\)。反过来，普通左 \(A\) 模沿 \(R\to A\) 取得 \(R\) 模结构，其作用双线性且平衡，故诱导模对象结构。两种同态条件也都成为 \(f(xv)=xf(v)\)。在集合的积结构中，这一过程给出幺半群在集合上的作用，而不是预先带加法的模。

\begin{proposition}\label{b4:prop:monads-as-monoidal-algebras}
设 \(\mathcal C\) 为小范畴。自函子范畴 \(\operatorname{Fun}(\mathcal C,\mathcal C)\) 以函子复合作为张量、以恒等函子作为单位，是严格幺半范畴。其代数对象恰好是 \(\mathcal C\) 上的单子。
\end{proposition}
\begin{proof}
对象的张量为 \(F\otimes G=F\circ G\)，自然变换的张量为水平复合。水平复合保持垂直复合与恒等自然变换，故确实给出双函子；这是自然变换交换律的内容。函子复合及自然变换的水平复合严格结合，恒等函子严格为单位，所以约束均可取恒等。

代数对象的对象现在是自函子 \(T\)。乘法 \(m:T\otimes T\to T\) 变成自然变换 \(\mu:T\circ T\Rightarrow T\)，单位 \(u:I\to T\) 变成 \(\eta:1_{\mathcal C}\Rightarrow T\)。左右张量 \(\mu\) 所得的分量分别是 \(\mu_{TX}:T^3X\to T^2X\)、\(T(\mu_X):T^3X\to T^2X\)，张量 \(\eta\) 所得的分量则是 \(\eta_{TX}\)、\(T(\eta_X)\)。因此代数对象公理逐分量写成
\[
\mu_X\circ\mu_{TX}=\mu_X\circ T(\mu_X),
\qquad
\mu_X\circ\eta_{TX}=1_{TX}=\mu_X\circ T(\eta_X),
\]
这恰好是\cref{b4:def:monad}的单子公理。反过来，每个单子给出这些相同的数据与等式，故两种结构完全对应。
\end{proof}

这一个例子中的张量是函子复合，通常不能交换两个因子。幺半结构能够容纳这种情形，因此代数对象不以辫状或对称结构为前提。

\subsection{内部复合与富集范畴}
\label{b4:subsec:monoidal-enrichment}

在线性代数中，同态不仅组成集合，还能相加；复合是双线性运算。因此把 Hom 作为对象，能够同时保留态射和态射之间的线性结构。内部 Hom 给出一种实现这一想法的方式。要把普通复合 \((g,f)\mapsto g\circ f\) 写成内部 Hom 之间的态射，可以先在输入上连续求值，再用转置确定所需态射；在模的例子中，这个求值过程就是 \((g,f,x)\mapsto g(f(x))\)。

\begin{proposition}\label{b4:prop:closed-internal-composition}
设 \(\mathcal C\) 为局部小闭幺半范畴。求值映射唯一确定内部复合
\[
\gamma_{X,Y,Z}:[Y,Z]\otimes[X,Y]\longrightarrow[X,Z]
\]
及单位 \(j_X:I\to[X,X]\)。其中 \(\gamma\) 是从 \(([Y,Z]\otimes[X,Y])\otimes X\) 到 \(Z\) 的下列态射的转置：
\[
\operatorname{ev}_{Y,Z}\circ
(1_{[Y,Z]}\otimes\operatorname{ev}_{X,Y})\circ
\alpha_{[Y,Z],[X,Y],X},
\]
而 \(j_X\) 是 \(\lambda_X:I\otimes X\to X\) 的转置。这些态射满足结合与左右单位等式。因此 \([X,X]\) 带有典范的代数对象结构。
\end{proposition}
\begin{proof}
转置的存在唯一性先给出 \(\gamma,j_X\)。对四个对象 \(W,X,Y,Z\)，比较
\[
\begin{gathered}
\gamma_{W,X,Z}\circ(\gamma_{X,Y,Z}\otimes1_{[W,X]}),\\
\gamma_{W,Y,Z}\circ(1_{[Y,Z]}\otimes\gamma_{W,X,Y})
\circ\alpha_{[Y,Z],[X,Y],[W,X]}.
\end{gathered}
\]
两者的定义域都是 \(([Y,Z]\otimes[X,Y])\otimes[W,X]\)。在模的例子中，给定 \(f:W\to X\)、\(g:X\to Y\)、\(h:Y\to Z\) 与 \(w\in W\)，两种复合都取值为 \(h(g(f(w)))\)。一般情形同样用连续求值比较：将这两个态射与 \(W\) 张量后求值，第一种复合先在 \(W\) 上求值到 \(X\)，再到 \(Y\)，最后到 \(Z\)；第二种也依次使用同样的三个求值。结合约束的五边形保证两种重括号到
\([Y,Z]\otimes([X,Y]\otimes([W,X]\otimes W))\) 相同，故所得态射相等。转置的唯一性因而给出结合等式。

将 \(j_Y\) 与 \([X,Y]\) 作内部复合，其在 \(X\) 上的求值仍是 \(\operatorname{ev}_{X,Y}\)，因为 \(j_Y\) 的转置是单位映射 \(I\otimes Y\to Y\)。三角等式删除这项单位，转置唯一性给出左单位律。将 \(j_X\) 放在另一侧，同样在 \(X\) 上求值得到原求值，给出右单位律。最后取三个对象都为 \(X\)，这些等式就是代数对象公理。
\end{proof}

在交换环 \(R\) 的模范畴中，内部复合为 \(g\otimes f\mapsto g\circ f\)，单位把 \(1\) 送到恒等同态。因此 \([P,P]=\operatorname{End}_R(P)\) 的代数对象结构就是通常的自同态代数。这也解释乘法次序为何按先 \(f\)、后 \(g\) 排列。
\symidx{\(\operatorname{End}_R(P)\)：模的自同态环}

普通局部小范畴为每对对象给出一个 Hom 集合。若把它换成 \(\mathcal V\) 中的对象，复合所用的集合笛卡尔积便换成 \(\mathcal V\) 的张量积，恒等元素也换成从单位对象 \(I\) 出发的态射。下面的定义保留结合与单位要求，以这样的 Hom 对象作为数据。

\begin{definition}\label{b4:def:enriched-category}
设 \(\mathcal V\) 为局部小幺半范畴，单位为 \(I\)。给定一个对象类，对每对对象 \(X,Y\) 给 \(\mathcal V\) 中的对象 \(\mathcal H(X,Y)\)，以及态射
\[
\gamma_{X,Y,Z}:\mathcal H(Y,Z)\otimes\mathcal H(X,Y)\longrightarrow\mathcal H(X,Z),
\qquad j_X:I\longrightarrow\mathcal H(X,X).
\]
若对任意 \(W,X,Y,Z\)，满足
\[
\begin{aligned}
\gamma_{W,X,Z}\circ(\gamma_{X,Y,Z}\otimes1)
={}&\gamma_{W,Y,Z}\circ(1\otimes\gamma_{W,X,Y})
\circ\alpha,\\
\gamma_{X,Y,Y}\circ(j_Y\otimes1)&=\lambda_{\mathcal H(X,Y)},\\
\gamma_{X,X,Y}\circ(1\otimes j_X)&=\rho_{\mathcal H(X,Y)},
\end{aligned}
\]
其中第一式的 \(\alpha\) 是 \(\mathcal H(Y,Z),\mathcal H(X,Y),\mathcal H(W,X)\) 的结合约束，恒等态射各在相应 Hom 对象上，则称这些数据为一个\term[fujifanchou]{富集范畴}[enriched category]，并称它富集于 \(\mathcal V\)。这里 Hom 对象不必是某个既有范畴的内部 Hom，而是所给的新数据。
\end{definition}
\symidx{\(\mathcal H(X,Y)\)：富集范畴中的态射对象}

当 \(\mathcal V=\mathbf{Ab}\)，张量为 \(\otimes_{\mathbb Z}\) 时，内部复合就是群同态 \(\mathcal H(Y,Z)\otimes\mathcal H(X,Y)\to\mathcal H(X,Z)\)，等价于双可加复合；\(\mathbb Z\to\mathcal H(X,X)\) 由恒等态射决定。因此这恰好恢复第四章预加性范畴的 Hom 群及复合。改用 \(k\)-向量空间，则恢复 \(k\) 线性范畴：复合对两个变量分别 \(k\) 线性。两者再要求零对象及有限双积存在，才得到加性范畴。

\ProseParagraphBreak

一般富集数据也给出普通范畴：取从 \(X\) 到 \(Y\) 的态射集合为 \(\operatorname{Hom}_{\mathcal V}(I,\mathcal H(X,Y))\)，即用从单位对象出发的态射来取 Hom 对象的“元素”。对 \(f:I\to\mathcal H(X,Y)\)、\(g:I\to\mathcal H(Y,Z)\)，张量 \(g\otimes f\) 的源是 \(I\otimes I\)，需要先接上单位同构的逆 \(\delta=\lambda_I^{-1}:I\to I\otimes I\)。在 Ab 或向量空间中，\(\delta\) 就是 \(1\mapsto1\otimes1\)。为与 \(\mathcal V\) 本身的态射复合区别，验证期间把新复合记为 \(\star\)：
\[
 g\star f\coloneqq\gamma_{X,Y,Z}\circ(g\otimes f)\circ\delta
 :I\longrightarrow\mathcal H(X,Z),
\]
恒等态射取为 \(j_X:I\to\mathcal H(X,X)\)。

\ProseParagraphBreak

若幺半结构严格，则 \(I\otimes I=I\)、\(\delta=1_I\)，新复合的结合律直接来自富集结合律。一般情形还需核对插入单位和改变括号的复合相容。单位约束满足 \(\lambda_I=\rho_I\)\cite[Corollary 2.2.5, p. 24]{EtingofEtAl2015TensorCategories}。将三角等式取 \(X=Y=I\) 并取逆，再前复合 \(\delta\)，得到
\[
 \alpha_{I,I,I}\circ(\delta\otimes1_I)\circ\delta
 =(1_I\otimes\delta)\circ\delta.
\]
结合 \(\alpha\) 的自然性与富集结合等式，对可相接的 \(f,g,h\) 得到 \((h\star g)\star f=h\star(g\star f)\)：两边插入三个单位时的差别恰由上式消去。若简记 \(H=\mathcal H(X,Y)\)，富集单位等式以及 \(\lambda,\rho\) 的自然性给出
\[
\begin{aligned}
 j_Y\star f&=\lambda_H\circ(1_I\otimes f)\circ\delta
 =f\circ\lambda_I\circ\delta=f,\\
 f\star j_X&=\rho_H\circ(f\otimes1_I)\circ\delta
 =f\circ\rho_I\circ\delta=f.
\end{aligned}
\]
这就验证了普通范畴的结合律与左右单位律；以后可把 \(\star\) 仍记为通常的复合符号 \(\circ\)。

\ProseParagraphBreak

在 Ab 和向量空间例子中，这正是取回 Hom 对象的元素。\chapref{b4:ch:20}的 DG 范畴把 Hom 对象改为复形，复合还必须与微分相容。

\ProseParagraphBreak
若 Hom 对象取为小范畴，张量取笛卡尔积，严格的复合与单位函子便给出严格2范畴。继续让 Hom 取值于前一层的严格范畴，就得到\secref{b4:sec:F03}的递归定义。
